%+TITLE: Mecanismo de Higgs
%+DESCRIPTION: Estudiamos cómo los bosones de gauge adquieren masa mediante la ruptura espontánea de simetría
% Fuente: transcripcion_teoria_clasica_campos_paginas_36-43/transcripcion_paginas_36-43.tex
% Páginas manuscritas 38--39.
% Se conserva la transcripción; los puntos pendientes de revisión se anotan abajo.
% Véase también transcripcion_teoria_clasica_campos_paginas_36-43/NOTAS_TRANSCRIPCION.md.
% Página 38, expansión del término cinético: en el segundo paréntesis aparece escrito \(-ig\,\partial^\mu\phi_2\), mientras que en el primero aparece \(+i\partial_\mu\phi_2\). También el término \(-\sqrt2igA^\mu\) parece carecer del factor \(a\). Ambos se preservaron.
% Página 38, término de mezcla: el manuscrito da \(+2\sqrt2g^2A_\mu\partial^\mu\phi_2\). Se conservó ese coeficiente aunque conviene revisarlo al editar.
% Página 38, transformación al gauge unitario: se escriben como igualdades exactas \(a+\phi_1'=a+\phi_1+\phi_2^2/(a+\phi_1)\) y \(\phi_1'=\phi_1+\phi_2^2/(a+\phi_1)\). Se mantuvieron sin introducir raíces ni aproximaciones.
% Página 39, expansión no abeliana: el último rótulo de la expansión parece decir «cuadráticos y cúbicos», aunque en contexto podría esperarse «cúbicos y cuárticos». Se transcribió lo visible.
\documentclass[11pt,a4paper]{article}

\usepackage[utf8]{inputenc}
\usepackage[T1]{fontenc}
\usepackage[spanish,es-nodecimaldot]{babel}
\usepackage{lmodern}
\usepackage{amsmath,amssymb,mathtools}
\usepackage{graphicx}
\usepackage{xcolor}
\usepackage{geometry}
\usepackage{enumitem}
\usepackage{microtype}

\geometry{margin=2.3cm}
\setlength{\parindent}{0pt}
\setlength{\parskip}{0.65em}
\setlist[itemize]{leftmargin=1.6em,itemsep=0.25em,topsep=0.25em}

\definecolor{title-red}{RGB}{165,25,20}
\definecolor{concept-green}{RGB}{20,125,45}
\newcommand{\concept}[1]{\textcolor{concept-green}{#1}}

\newcommand{\ket}[1]{\left|#1\right\rangle}
\newcommand{\bra}[1]{\left\langle#1\right|}
\newcommand{\braket}[2]{\left\langle #1\middle|#2\right\rangle}
\newcommand{\hc}{\mathrm{h.c.}}
\newcommand{\Tr}{\operatorname{Tr}}
\newcommand{\Diffs}{\mathrm{Diffs}}
\newcommand{\slashedpartial}{\not\!\partial}

\begin{document}

% -----------------------------------------------------------------------------
% Página manuscrita 38
% -----------------------------------------------------------------------------
Vimos que la ruptura espontánea de una simetría global garantiza la presencia
de partículas sin masa. Así mismo, la simetría de gauge garantiza la presencia
de bosones de gauge sin masa. Pero cuando la simetría de gauge se rompe
espontáneamente lo que ocurre es que los bosones de gauge adquieren masa. Esto
se llama mecanismo de Higgs.

Empecemos por estudiar la ruptura espontánea de una simetría de gauge $U(1)$.
Para eso consideramos el lagrangiano
\[
  \mathcal L
  =-D_\mu\phi^*D^\mu\phi
   -m^2\phi^*\phi
   -\lambda(\phi^*\phi)^2
   -\frac14F_{\mu\nu}F^{\mu\nu}.
\]

Donde
\[
  D_\mu\phi=\partial_\mu\phi+igA_\mu\phi,
  \qquad
  F_{\mu\nu}=\partial_\mu A_\nu-\partial_\nu A_\mu.
\]

Si $m^2<0$, el mínimo de potencial es en
\[
  \phi^*\phi=-\frac{m^2}{2\lambda}=a^2.
\]
Sin pérdida de generalidad, tomamos el mínimo en $\phi_0=a$ (sin parte
imaginaria). Entonces parametrizamos el campo como
\[
  \phi=a+\frac1{\sqrt2}(\phi_1+i\phi_2),
\]
obteniendo
\begin{align*}
  \mathcal L
  ={}&-\frac12
  \left(
    \partial_\mu\phi_1+i\partial_\mu\phi_2
    +\sqrt2\,igaA_\mu+\text{cuadráticos}
  \right)
  \left(
    \partial^\mu\phi_1-ig\partial^\mu\phi_2
    -\sqrt2\,igA^\mu+\text{cuadráticos}
  \right)\\
  &+\lambda a^2\left((\sqrt2a+\phi_1)^2+\phi_2^2\right)
  -\frac\lambda4\left((\sqrt2a+\phi_1)^2+\phi_2^2\right)^2
  -\frac14F_{\mu\nu}F^{\mu\nu}\\
  ={}&-\frac12\partial_\mu\phi_1\partial^\mu\phi_1
  -2\lambda a^2\phi_1^2
  -\frac12\partial_\mu\phi_2\partial^\mu\phi_2
  -\frac14F_{\mu\nu}F^{\mu\nu}
  +2g^2a^2A_\mu A^\mu
  +2\sqrt2g^2A_\mu\partial^\mu\phi_2\\
  &+\text{términos cúbicos y cuárticos}.
\end{align*}

Como antes, tenemos un campo sin masa $\phi_2$ relacionado con la dirección
rota. Pero tenemos un par de novedades. Aparece un término de masa para el
bosón de gauge y aparece un término cuadrático que involucra $\phi_2$ y el bosón
de gauge $A_\mu$. Ese término cuadrático se llama término de mezcla porque en la
teoría cuántica implica que $\phi_2$ se puede transformar en una componente de
$A_\mu$ y vice versa.

En realidad tenemos demasiados campos. La ambigüedad de gauge hace que podamos
eliminar algunos. Esto es análogo a eliminar coordenadas escogiendo
inteligentemente el marco de referencia. Para ver cómo hagamos
\[
  \phi\longmapsto\phi'=e^{i\alpha(x)}\phi,
\]
por lo que
\[
  a+\phi_1'=(a+\phi_1)\cos\alpha(x)-\phi_2\sin\alpha(x),
\]
\[
  \phi_2'=\phi_2\cos\alpha(x)+(a+\phi_1)\sin\alpha(x).
\]

Basta escoger
\[
  \tan\alpha(x)=-\frac{\phi_2}{a+\phi_1},
\]
que podemos hacer porque $\alpha$ depende de $x$. De este modo
\[
  \phi_2'=0,
  \qquad
  a+\phi_1'=a+\phi_1+\frac{\phi_2^2}{a+\phi_1},
  \qquad
  \phi_1'=\phi_1+\frac{\phi_2^2}{a+\phi_1}.
\]

En este gauge (llamado gauge \concept{unitario}) el lagrangiano toma una forma
sencilla:
\[
  \mathcal L
  =-\frac12\partial_\mu\phi_1\partial^\mu\phi_1
   -2a^2\lambda\phi_1^2
   -\frac14F_{\mu\nu}F^{\mu\nu}
   +2g^2a^2A_\mu A^\mu
   +\text{términos cúbicos y cuárticos}.
\]

% -----------------------------------------------------------------------------
% Página manuscrita 39
% -----------------------------------------------------------------------------
En este lagrangiano desaparece el bosón de Goldstone, pero aparece un término de
masa para el bosón de gauge. Un campo vectorial sin masa tiene dos grados de
libertad, mientras uno con masa tiene tres (tarea). Decimos entonces que
``el bosón de gauge se come al bosón de Goldstone para adquirir masa'' (entre
comillas).

Ahora hagamos lo mismo para una teoría no abeliana. Veremos que son los bosones
de gauge correspondientes a las direcciones rotas los que adquieren masa.

Para simplificar la notación mencionamos un hecho de la teoría de grupos. Para
cualquier grupo compacto (por ejemplo $SU(n)$ o $SO(n)$) se pueden escoger los
generadores tales que
\[
  \Tr(J^iJ^j)=\frac12\delta^{ij}
  \qquad\Longrightarrow\qquad
  -\frac12\Tr\left[
    (F_{\mu\nu}^iJ^i)(F^{j\mu\nu}J^j)
  \right]
  =-\frac14F_{\mu\nu}^iF^{i\mu\nu}.
\]

Nuestro lagrangiano de partida es (con $\vec\phi=(\phi_1,\phi_2,\phi_3)$)
\[
  \mathcal L
  =-\frac12(D_\mu\phi_i)(D^\mu\phi_i)
   -\frac{m^2}{2}\phi_i\phi_i
   -\lambda(\phi_i\phi_i)^2
   -\frac14F_{\mu\nu}^iF^{i\mu\nu},
\]
con
\[
  D_\mu\phi_i=\partial_\mu\phi_i+g\epsilon_{ijk}A_\mu^j\phi_k,
\]
\[
  F_{\mu\nu}^i
  =\partial_\mu A_\nu^i-\partial_\nu A_\mu^i
   +g\epsilon_{ijk}A_\mu^jA_\nu^k,
\]
donde hemos usado el álgebra de Lie
\[
  [J_i,J_j]=i\epsilon_{ijk}J_k.
\]
Del caso global sabemos que si $m^2<0$ el mínimo del potencial está en
\[
  \phi_{0i}\phi_{0i}=-\frac{m^2}{4\lambda}=a^2,
  \qquad
  \text{sin pérdida de generalidad}\qquad
  \vec\phi_0=a\hat e_3.
\]

Parametrizamos el campo
\[
  \vec\phi=(\phi_1,\phi_2,\chi+a)
\]
y el lagrangiano cuadrático queda
\begin{align*}
  \mathcal L={}&-\frac12\Bigl[
    (\partial_\mu\phi_1+g\epsilon_{123}A_\mu^2a)^2
    +(\partial_\mu\phi_2+g\epsilon_{213}A_\mu^1a)^2
    +(\partial_\mu\chi)^2
    +\text{cúbicos y cuárticos}
  \Bigr]\\
  &+2\lambda a^2\left(\phi_1^2+\phi_2^2+(\chi+a)^2\right)
   -\lambda\left(\phi_1^2+\phi_2^2+(\chi+a)^2\right)^2
   -\frac14F_{\mu\nu}^iF^{i\mu\nu}\\
  ={}&-\frac12\partial_\mu\chi\partial^\mu\chi
   -2\lambda a^2\chi^2
   -\frac12\partial_\mu\phi_1\partial^\mu\phi_1
   -\frac12\partial_\mu\phi_2\partial^\mu\phi_2\\
  &-ga\,\partial_\mu\phi_1A^{2\mu}
   +ga\,\partial_\mu\phi_2A^{1\mu}
   -\frac14F_{\mu\nu}^iF^{i\mu\nu}
   -\frac12g^2a^2A_\mu^1A^{1\mu}
   -\frac12g^2a^2A_\mu^2A^{2\mu}\\
  &+\text{cuadráticos y cúbicos}.
\end{align*}

Tenemos términos de masa para $A_\mu^1,A_\mu^2$ que corresponden a los
generadores rotos $J^1,J^2$, mientras que $A_\mu^3$ sigue sin masa. Aparecen dos
bosones de Goldstone sin masa pero mezclados con $A_\mu^1$ y $A_\mu^2$. Estos se
pueden eliminar de nuevo realizando una transformación de gauge al gauge
unitario, reorientando el campo punto por punto para que apunte en la dirección
$\hat e_3$:
\[
  \vec\phi=(a+\chi)\hat e_3.
\]

Entonces el lagrangiano cuadrático queda
\[
  \mathcal L
  =-\frac12\partial_\mu\chi\partial^\mu\chi
   -2\lambda a^2\chi^2
   -\frac14F_{\mu\nu}^iF^{i\mu\nu}
   -\frac12g^2a^2A_\mu^1A^{1\mu}
   -\frac12g^2a^2A_\mu^2A^{2\mu}
   +\cdots.
\]

``Los bosones de gauge en las direcciones rotas se comen a los bosones de
Goldstone para adquirir masa.''

Al campo $\vec\phi$ se lo llama un \concept{campo de Higgs}.

\end{document}
